\documentclass[10pt,a4paper]{amsart}
\usepackage[margin=19mm]{geometry}
\usepackage{amsmath,amssymb,mathtools}
\usepackage[scr=rsfs,cal=euler]{mathalfa}
\usepackage{xfrac}
\usepackage[unicode,hidelinks,bookmarks=false]{hyperref}
\newcommand{\ie}{i.e.\ }
\newcommand{\cf}{cf.\ }
\newcommand{\resp}{respectively\ }
\newcommand{\Subsection}{\subsection}
\providecommand{\nobreakdash}{\nobreak\hbox{-}}
\setlength{\emergencystretch}{2em}
\multlinegap0pt
\usepackage{bm}
\usepackage{bbm}
\usepackage{dsfont}
\usepackage{xspace}
\usepackage{cancel}
\usepackage{mathtools}
\usepackage{amsthm}
\makeatletter
\@ifundefined{lemma}{\newtheorem{lemma}{Lemma}}{}
\makeatother
\title[Asymptotics and zero behaviour of geometric polynomials]{Asymptotics and zero behaviour of geometric polynomials}
\author{M. Bello-Hern\'andez, M. Benito, \'O. Ciaurri,, E. Fern\'andez}
\date{Source: arXiv:2602.15718v2 (CC BY 4.0)}
\begin{document}
\maketitle

\noindent\emph{Source and licence.} Asymptotics and zero behaviour of geometric polynomials. Version arXiv:2602.15718v2 (CC BY 4.0), distributed under the Creative Commons Attribution 4.0 International licence, as recorded in the arXiv licence metadata for this exact version and shown on the abstract page https://arxiv.org/abs/2602.15718v2. Full rights notice: licenses/arxiv-2602.15718-CC-BY-4.0.txt.\par\smallskip

\noindent\textit{Editorial scope.} Counted unit: the Lemma whose source label is \texttt{LemCocienteCons} in the article (source0.tex lines 322--332, source label \texttt{LemCocienteCons}), delivered together with the article's own complete original proof of it (source0.tex lines 334--367, one \texttt{proof} environment of 34 lines). No mathematics is written, repaired or re-derived: statement and proof are the article's own text line for line. Required context retained uncounted: lemZer (lines 233--259). Every definition, display and label of those lines is retained unchanged. Omitted from this package and not redistributed: 1--232, 260--321, 333--333, 368--974. Nothing inside a retained excerpt is omitted or altered: every retained body is delivered exactly as the article's lines are written, so no omission receipt of any kind accompanies this package. Citations: the retained excerpts carry no citation command at all, so no bibliography entry of the article is transcribed and no reference is redistributed. Reference adaptation: none was needed and none was made; every reference inside the delivered unit resolves inside it, so no unresolved reference or citation is produced. Printed slips kept literal: no printed word, symbol, hypothesis or number of the counted statement or of its proof was altered. Layout and class: the article's own class is \texttt{article} with packages including amsmath, amssymb, amsthm, inputenc, babel, graphicx, color, hyperref, multicol, subcaption; the wrapper is the lane's pinned \texttt{amsart} editorial wrapper at 10pt on A4, which restates the theorem-like environments the retained text needs and adds the article's own macro definitions that the retained text needs (none), with extra wrapper packages (bm, bbm, dsfont, xspace, cancel, mathtools). The delivered numbering is the unit's own. Assets: no retained span contains a figure environment, a graphics-inclusion command, a PDF-page inclusion, a table or a TikZ picture; no image file is copied into this package, the package declares no asset dependency and no image file is redistributed with it. Licence: version arXiv:2602.15718v2 of this article is distributed under the Creative Commons Attribution 4.0 International licence, as recorded in the arXiv licence metadata for this exact version and shown on the abstract page https://arxiv.org/abs/2602.15718v2. A full rights notice is delivered at licenses/arxiv-2602.15718-CC-BY-4.0.txt. Characters: every retained excerpt is pure ASCII, so the delivered engine, XeLaTeX with its default Latin Modern text font, reports no missing character.
\begin{lemma}\label{lemZer} 
\begin{enumerate}
\item
The following formula holds
\begin{equation}
\label{delDifEc}
P_n(z)=z\frac{d((1+z)P_{n-1}(z))}{dz}=z\,P_{n-1}(z)+z(1+z)\,P_{n-1}'(z),
\end{equation}
and $P_n$ has simple zeros in $(-1,0]$. Moreover the zeros of $P_n$ and $P_{n-1}$  different from $z=0$ interlace.

\item We have
\begin{equation}
\label{simtRel}
(1+z)P_n(z)=(-1)^nzP_n(-1-z),\quad n\ge 1.
\end{equation}
The zeros of $P_n$ different from $z=0$ are symmetric with respect to $z=-1/2$. If $n$ is even, $P_n$ has  a zero at $z=-1/2$. In addition
\begin{equation}
\label{RelPolDer}
(1+z)P_n'(z)+(-1)^nzP_n'(-1-z)=\frac{P_n(z)}{z}.	
\end{equation}
If $P_n(z)\ne 0$ and $z\ne -1$, it holds
\begin{equation}
\label{RelPolDer2}
\frac{P_n'(z)}{P_n(z)}+\frac{P_n'(-1-z)}{P_n(-1-z)}=\frac{1}{z(1+z)}.
\end{equation}
\end{enumerate}
\end{lemma}

\begin{lemma}\label{LemCocienteCons} There exist $n-1$ positive numbers $\lambda_{k,n}$ such that 
\begin{equation}
\label{descFracSimples}
n\frac{P_{n-1}(z)}{P_n(z)}=\sum_{k=1}^{n-1}\frac{\lambda_{k,n}}{z-z_{k,n}^*}
\end{equation}
and $\sum_{k=1}^{n-1}\lambda_{k,n-1}=1$. Moreover the sequence of functions $\left(n\frac{P_{n-1}(z)}{P_n(z)}\right)_{n\in\mathbb{N}}$ is uniformly bounded in compact subsets of  $\mathbb{C}\setminus[-1,0]$; i.e. if $K\subset \overline{\mathbb{C}}\setminus[-1,0]$ is a compact set, then there exists a constant $C=C(K)>0$ such that for all  $z\in K$ we have
\[
\left|n\frac{P_{n-1}(z)}{P_n(z)}\right|\le C.
\]

\end{lemma}

\begin{proof} According to  Lemma \ref{lemZer} the zeros of $P_n$ are simple and lie in $(-1,0]$. Let
\[
Q_k(z)=\frac{P_k(z)}{z},\quad k=1,2,\ldots
\]
The zeros of $Q_n$ are the $n-1$ numbers $(z_{k,n}^*)$. By the Lagrange  interpolation formula we have
\[
nQ_{n-1}(z)=\sum_{k=1}^{n-1}nQ_{n-1}(z_{k,n}^*)\frac{Q_n(z)}{Q_n'(z_{k,n}^*)(z-z_{k,n}^*)}.
\]
Thus, the equality \eqref{descFracSimples} follows with
\[
\lambda_{k,n}=\frac{nQ_{n-1}(z_{k,n}^*)}{Q_n'(z_{k,n}^*)}.
\]
The polynomials $Q_n$ y $Q_{n-1}$ have degree exactly $n-1$ and $n-2$, respectively, with leading coefficients $n!$ and $(n-1)!$ Multiplying  \eqref{descFracSimples} by $z$ and taking limits as $z\to\infty$, we obtain
\[
\sum_{k=1}^{n-1}\lambda_{k,n}=1.
\]
According to  Lemma \ref{lemZer}, the zeros of $Q_n$ and $Q_{n-1}$  interlace in $(-1,0)$. From Rolle's theorem $Q_n$ and $Q_n'$ have also interlacing zeros, so 
\[
Q_n'(z_{k,n}^*)Q_n'(z_{k-1,n}^*)<0,
\]
the same happen with $Q_{n-1}$. As $Q_{n-1}$ and $Q_{n}$ have positive leading coefficient, we have
$Q_n'(z_{k,n}^*)Q_n(z_{k,n}^*)>0$ and hence $\lambda_{k,n}>0$, $k=1,2,\ldots,n-1.$

If $K\subset \mathbb{C}\setminus [-1,0]$ is a compact set, 
\[
d(K,[-1,0])=\inf\{|z-x|; x\in[-1,0],\, z\in K\}>0.
\]
From \eqref{descFracSimples} it follows
\[
\left|n\frac{P_{n-1}(z)}{P_n(z)}\right|\le \frac{1}{d(K,[-1,0])}.\qedhere
\]


\end{proof}

\end{document}
